% objc-interop.tex — the Objective-C an iOS codebase still contains (messaging, nil,
% id, properties, blocks, categories, KVC/KVO, swizzling) + Swift <-> ObjC
% interop (@objc, dynamic, NSObject, headers, nullability, what can't cross).
% Sources (checked 2026-10-06):
%   developer.apple.com/library/archive Objective-C Runtime Programming Guide (messaging,
%     dynamic method resolution, forwarding) · Key-Value Observing Programming Guide
%     (isa-swizzling) · Blocks Programming Topics (stack blocks copied to the heap)
%   developer.apple.com/documentation/swift/importing-objective-c-into-swift ·
%     /using-imported-c-enums-in-swift (NS_ENUM / NS_OPTIONS) · /improving-objective-c-api-
%     declarations-for-swift (NS_REFINED_FOR_SWIFT) · /about-imported-cocoa-error-parameters
%     (NSError ** <-> throws) · documentation/objectivec/nsobject/load + initialize
%   swift-evolution SE-0160 (no @objc inference, Swift 4.0) — checked against the index
% @labels: area=swift kind=concept platform=apple topic=language discussion=147
% @tags: objective-c, objc-msgsend, message-forwarding, objc-dynamic, nsobject, bridging-header, nullability, kvo, kvc, method-swizzling, categories, blocks
\documentclass[8pt]{extarticle}
\usepackage{printup-sheet}

\lstdefinelanguage{SwiftSheet}{
  morekeywords={protocol,class,final,struct,enum,func,var,let,weak,init,
    if,else,return,guard,self,nil,true,false,dynamic,override,private,
    NSObject,Double,String,Int,Any},
  sensitive=true, morecomment=[l]{//}, morecomment=[s]{/*}{*/}, morestring=[b]"}
\lstdefinelanguage{ObjCSheet}{
  alsoletter={@_},
  morekeywords={@interface,@implementation,@end,@property,nonatomic,strong,
    weak,copy,assign,readonly,id,nil,self,void,return,if,YES,NO,BOOL,
    __weak,__strong,__block,typeof,nullable,nonnull,instancetype},
  sensitive=true, morecomment=[l]{//}, morestring=[b]"}

\tikzset{
  sb/.style={box, font=\scriptsize, inner sep=1.5pt, minimum height=4.6mm},
  ok/.style={sb, draw=sheetGreen, fill=sheetGreen!10},
  bad/.style={sb, draw=sheetRed, fill=sheetRed!8},
  q/.style={sb, draw=sheetOrange, fill=sheetOrange!10},
  lbl/.style={font=\tiny, text=black!75, inner sep=1pt, align=center},
  ttl/.style={font=\bfseries\small, anchor=west},
}

\begin{document}

\sheettitle{Objective-C essentials + Swift $\leftrightarrow$ ObjC interop}{swift · folio}

\oneliner{An ObjC call is a \textbf{message}: \texttt{[obj sel]} becomes
\texttt{objc\_msgSend(obj, @selector(sel))}, resolved \textbf{at runtime} in the
receiver's class — hence nil messaging, categories, KVO, swizzling. Swift enters
that world via \texttt{@objc} (\emph{visible} to the runtime) and \texttt{dynamic}
(\emph{dispatched} through it), on \texttt{NSObject}-rooted classes.}

\vspace{2pt}
\noindent\begin{tikzpicture}[sheet]
  \node[ttl] at (-0.1,2.45) {\texttt{[obj tap:x]} $\to$ \texttt{objc\_msgSend(obj, @selector(tap:), x)}};
  \node[q] (nil) at (0.7,1.6) {obj nil?};
  \node[ok, text width=13mm] (ret0) at (0.7,0.45) {return 0 / nil /\\\texttt{NO} — no crash};
  \node[sb] (cache) at (2.45,1.6) {isa $\to$ class\\method cache};
  \node[sb] (list) at (4.35,1.6) {method list\\(+ categories)};
  \node[sb] (sup) at (6.25,1.6) {superclass\\chain};
  \node[ok] (imp) at (4.35,0.45) {IMP found $\to$ tail-call\\\texttt{imp(obj, \_cmd, x)}};
  \node[q] (res) at (8.2,1.6) {\texttt{+resolve}\\\texttt{InstanceMethod:}};
  \node[q] (fwd) at (10.2,1.6) {\texttt{forwarding}\\\texttt{TargetForSelector:}};
  \node[q] (inv) at (12.2,1.6) {\texttt{forward}\\\texttt{Invocation:}};
  \node[bad] (crash) at (12.2,0.45) {\texttt{doesNotRecognize}\\\texttt{Selector:} $\to$ crash};
  \draw[flow] (nil) -- node[lbl, left]{yes} (ret0);
  \draw[flow] (nil) -- node[lbl, above]{no} (cache);
  \draw[flow] (cache) -- node[lbl, above]{miss} (list);
  \draw[flow] (list) -- node[lbl, above]{miss} (sup);
  \draw[flow] (sup) -- node[lbl, above]{miss} (res);
  \draw[flow] (res) -- (fwd); \draw[flow] (fwd) -- (inv);
  \draw[flow] (inv) -- node[lbl, right]{no} (crash);
  \draw[flow, draw=sheetGreen] (cache.south) -- node[lbl, left]{hit} (imp.north west);
  \draw[flow, draw=sheetGreen] (list) -- (imp);
  \draw[flow, draw=sheetGreen] (sup.south) -- (imp.north east);
  \node[lbl, text=sheetBrown] at (9.5,0.45) {add a method → other target\\→ full forward (\texttt{NSInvocation})};
  \node[lbl, text=sheetRed] at (6.55,0.45) {``unrecognized selector\\sent to instance''};
  \draw[sheetGrey!40] (13.75,2.6) -- (13.75,0.1);
  \node[ttl] at (13.8,2.45) {two headers, two directions};
  \node[sb, fill=sheetGreen!10, draw=sheetGreen] (sw) at (14.4,1.35) {Swift};
  \node[sb, fill=sheetOrange!10, draw=sheetOrange] (oc) at (16.3,1.35) {ObjC};
  \draw[flow] (oc.north) to[out=150, in=30] node[lbl, above]{\texttt{T-Bridging-Header.h}} (sw.north);
  \draw[flow] (sw.south) to[out=-30, in=-150] node[lbl, below]{\texttt{Module-Swift.h} (generated)} (oc.south);
\end{tikzpicture}

\begin{multicols}{2}
\raggedright

\section{ObjC}
\begin{itemize}
  \item \texttt{SEL} = interned name; \texttt{IMP} = C fn
        \texttt{(id self, SEL \_cmd, ...)}. Cached lookup, never inlined.
        \textbf{nil} receiver returns zero — a silent bug, not a crash.
        \texttt{id} = any object, unchecked (Swift: \texttt{Any}).
  \item \textbf{Properties:} default \texttt{atomic} (accessor lock only, not
        thread safety). \texttt{strong} own · \texttt{weak} zeroing ·
        \texttt{copy} for \texttt{NSString}/blocks (caller may pass a mutable
        one) · \texttt{assign} scalars only (an object dangles; delegates are
        \texttt{weak}) · \texttt{unsafe\_unretained} not zeroed.
  \item \textbf{Blocks} capture locals \textbf{by value} at creation and
        \textbf{retain} captured objects; \texttt{\_\_block} = by reference,
        mutable. Born on the stack, ARC copies to heap when stored.
        \texttt{self$\to$block$\to$self} cycle: weak/strong dance.
  \item \textbf{Categories} add methods (not ivars — use associated objects)
        to any class. Same selector twice = undefined winner $\Rightarrow$
        prefix. \texttt{@interface C ()} = private extension in \texttt{.m}.
  \item \textbf{KVC} \texttt{valueForKeyPath:@"a.b"}, \texttt{@sum.x};
        bad key $\to$ \texttt{NSUnknownKeyException}. \textbf{KVO}
        isa-swizzles the object to hidden \texttt{NSKVONotifying\_C} whose
        setter wraps \texttt{will/didChangeValueForKey:}; fires on the
        \emph{mutating} thread.
  \item \textbf{Swizzle} with \texttt{method\_exchangeImplementations} in
        \texttt{+load} + \texttt{dispatch\_once} (\texttt{+load}: pre-main,
        per class and category; \texttt{+initialize}: lazy, may run twice).
\end{itemize}

\section{Example}
\begin{lstlisting}[language=ObjCSheet]
@property (nonatomic, copy) NSString *title;
@property (nonatomic, weak, nullable) id<PlayerDelegate> delegate;
@property (nonatomic, copy) void (^onEnd)(void);
__weak typeof(self) weakSelf = self;
self.onEnd = ^{
  __strong typeof(weakSelf) s = weakSelf;   // alive for the body
  if (!s) return;
  [s.delegate playerDidEnd:s]; };           // nil delegate: no-op
\end{lstlisting}
\begin{lstlisting}[language=SwiftSheet]
final class Model: NSObject {
  @objc dynamic var progress = 0.0     // KVO works
  @objc var name = ""                  // visible, NOT KVO-safe
  @objc(loadWithID:) func load(id: Int) {} }
let token = model.observe(\.progress, options: [.new]) { m, c in
  print(c.newValue as Any) }   // hold token; its deinit removes
\end{lstlisting}

\section{ObjC $\to$ Swift import}
{\footnotesize
\begin{tabular}{@{}ll@{\hspace{4mm}}ll@{}}
\texttt{id} & \texttt{Any} & \texttt{Class} & \texttt{AnyClass} \\
\texttt{instancetype} & \texttt{Self} & \texttt{SEL} & \texttt{Selector} \\
\texttt{BOOL} & \texttt{Bool} & block & \texttt{@convention(block)} closure \\
\texttt{NSInteger} & \texttt{Int} & category & \texttt{extension} \\
\end{tabular}}

\columnbreak

\section{Interop}
\begin{itemize}
  \item \texttt{@objc} exposes (selector, \texttt{-Swift.h}) but may still be
        vtable-called; \texttt{dynamic} forces \texttt{objc\_msgSend}
        $\Rightarrow$ KVO, swizzling. \texttt{@objcMembers} = all members.
        Swift 4 (SE-0160) stopped inferring \texttt{@objc}.
  \item \texttt{@objc} class must inherit \texttt{NSObject};
        \texttt{@objc protocol} is class-only, allows \texttt{optional}.
  \item Bridging header: app/test targets (frameworks use umbrella header).
        \texttt{-Swift.h}: \texttt{internal}+ in an app, \texttt{public} in a
        framework; \texttt{\#import} in \texttt{.m}, \texttt{@class} in \texttt{.h}.
  \item \texttt{nonnull}$\to$\texttt{T} · \texttt{nullable}$\to$\texttt{T?} ·
        unannotated$\to$\texttt{T!}; use \texttt{NS\_ASSUME\_NONNULL\_BEGIN}.
  \item \texttt{NS\_SWIFT\_NAME} renames · \texttt{NS\_REFINED\_FOR\_SWIFT}
        $\to$ \texttt{\_\_name} · \texttt{NS\_ENUM}$\to$enum (\texttt{@unknown
        default}) · \texttt{NS\_OPTIONS}$\to$\texttt{OptionSet} ·
        \texttt{NSArray<NSString *>}$\to$\texttt{[String]} (else \texttt{[Any]}) ·
        \texttt{NSError **}$\leftrightarrow$\texttt{throws}.
  \item \textbf{Can't cross:} structs (bar bridged \texttt{String}/\texttt{Data}/%
        \texttt{URL}…), enums with payloads (only \texttt{@objc enum E: Int}),
        generics, tuples, \texttt{Int?} (use \texttt{NSNumber?}), PATs /
        Swift-only protocols. Default args cross \emph{without} defaults.
\end{itemize}

\section{Traps}
\begin{itemize}
  \trap{\texttt{@objc} alone: Swift callers use the vtable and bypass KVO's
        setter or a swizzled IMP — observation silently never fires.}
  \trap{Swizzling: global, order-dependent, \texttt{\_cmd} changes; on an
        \emph{inherited} method it swaps the superclass's
        (\texttt{class\_addMethod} first); breaks across iOS versions.}
  \trap{\texttt{[Int]}$\to$\texttt{NSArray} boxes each element in
        \texttt{NSNumber}; \texttt{String}$\leftrightarrow$\texttt{NSString}
        may transcode — costly in a loop.}
  \trap{Overloads collide as one selector — \texttt{@objc(name:)}. No
        \texttt{@objc} on an IB action $\to$ ``unrecognized selector''.}
  \trap{Header says \texttt{nonnull}, ObjC returns nil $\to$ Swift crashes
        on use. \texttt{@Published} is Combine, not KVO.}
\end{itemize}

\section{Remember}
\textbf{``@objc shows, dynamic sends.''} nil eats messages; categories add
methods not storage; KVO = a secret subclass.


\end{multicols}

\noindent{\footnotesize\color{sheetGrey}\textit{Related:} \texttt{arc-ownership}
· \texttt{observation-and-combine} (KVO vs Combine) · \texttt{protocols-generics}
(dispatch) · \texttt{swift-unsafe-and-reflection} (\texttt{Unmanaged}) ·
\texttt{swift-runtime-and-compiler}}

\end{document}
